\section*{Capítulo \ref{ch:b1:crystals-metals} --- Cristales I: el cristal perfecto y los metales}

\begin{solution}{exo:b1:crystals-metals:1}
Cúbica simple: $8 \times \frac18 = 1$ átomo, coordinación 6. BCC: $8 \times
\frac18 + 1 = 2$, coordinación 8. FCC: $8 \times \frac18 + 6 \times
\frac12 = 4$, coordinación 12.
\end{solution}

\begin{solution}{exo:b1:crystals-metals:2}
Los átomos se tocan a lo largo de una arista: $a = 2R$, $Z = 1$, $C = \frac43\pi R^3/(2R)^3
= \pi/6 \approx 0.52$, menos que en la BCC (0.68) y la FCC (0.74).
\end{solution}

\begin{solution}{exo:b1:crystals-metals:3}
$R = a\sqrt2/4 = 405.0 \times 0.3536 = \qty{143.2}{pm}$;
$\rho = 4 \times 26.98/(6.022 \times 10^{23} \times (4.050 \times
10^{-8})^3) = \qty{2.70}{g/cm^3}$.
\end{solution}

\begin{solution}{exo:b1:crystals-metals:4}
$R = a\sqrt3/4 = \qty{185.8}{pm}$; $\rho = 2 \times 22.99/(6.022 \times
10^{23} \times (4.291 \times 10^{-8})^3) = \qty{0.966}{g/cm^3}$, de
acuerdo con el valor medido, \qty{0.97}{g/cm^3}.
\end{solution}

\begin{solution}{exo:b1:crystals-metals:5}
$a^3 = 4M/(\rho N_A) = 4 \times 107.87/(10.50 \times 6.022 \times 10^{23}) =
\qty{6.824e-23}{cm^3}$, de modo que $a = \qty{4.086e-8}{cm} = \qty{408.6}{pm}$ y
$R = a\sqrt2/4 = \qty{144.5}{pm}$.
\end{solution}

\begin{solution}{exo:b1:crystals-metals:6}
La demostración es la de la \cref{prop:b1:crystals-metals:hcp}. Magnesio:
$c/a = 521.0/320.9 = 1.624$. El prisma, de volumen
$\frac{3\sqrt3}{2}a^2c$, contiene 6 átomos, de modo que $\rho = 6M/(N_A \cdot
\frac{3\sqrt3}{2}a^2c) = 6 \times 24.31/(6.022 \times 10^{23} \times
2.598 \times (3.209 \times 10^{-8})^2 \times 5.210 \times 10^{-8}) =
\qty{1.74}{g/cm^3}$.
\end{solution}

\begin{solution}{exo:b1:crystals-metals:7}
Átomos: $(0,0,0)$, $(\frac a2,\frac a2,0)$, $(\frac a2,0,\frac a2)$,
$(0,\frac a2,\frac a2)$. \omterm{def:b1:crystals-metals:site}{Huecos octaédricos}: $(\frac a2,\frac a2,\frac a2)$,
$(\frac a2,0,0)$, $(0,\frac a2,0)$, $(0,0,\frac a2)$. \omterm{def:b1:crystals-metals:site}{Huecos tetraédricos}:
los 8 puntos cuyas coordenadas valen cada una $\frac a4$ o
$\frac{3a}{4}$. Un \omterm{def:b1:crystals-metals:site}{hueco octaédrico} y dos tetraédricos por átomo.
\end{solution}

\begin{solution}{exo:b1:crystals-metals:8}
$R(\ce{Cu}) = 361.5\sqrt2/4 = \qty{127.8}{pm}$, $R(\ce{Ni}) = 352.4\sqrt2/4
= \qty{124.6}{pm}$: radios que difieren menos de un 3\,\% y la misma
estructura, de modo que cualquiera de los dos átomos puede ocupar el
lugar del otro en la \omterm{def:b1:crystals-metals:lattice}{red}. El hidrógeno, mucho más pequeño, solo puede
alojarse en \omterm{def:b1:crystals-metals:site}{huecos intersticiales}: forma una \omterm{def:b1:crystals-metals:alloy}{aleación intersticial}.
\end{solution}

\begin{solution}{exo:b1:crystals-metals:9}
$R = 316.5\sqrt3/4 = \qty{137.0}{pm}$; $\rho = 2 \times 183.84/(6.022 \times
10^{23} \times (3.165 \times 10^{-8})^3) = \qty{19.26}{g/cm^3}$; átomos por
\unit{cm^3}: $2/a^3 = 2/(3.165 \times 10^{-8})^3 = \num{6.31e22}$.
\end{solution}

\begin{solution}{exo:b1:crystals-metals:10}
El centro de cara $(\frac a2,\frac a2,0)$ está a $a/2$ de los centros de
las celdas de arriba y de abajo: $r_O = \frac a2 - R = \frac a2 - \frac{a\sqrt3}{4}
= \frac{2-\sqrt3}{4}a$. El punto $(\frac a2,\frac a4,0)$ está a
$\sqrt{a^2/4 + a^2/16} = \frac{a\sqrt5}{4}$ de los vértices $(0,0,0)$ y
$(a,0,0)$ y de los centros $(\frac a2,\frac a2,\pm\frac a2)$:
$r_T = \frac{\sqrt5 - \sqrt3}{4}a$. Con $a = 4R/\sqrt3$:
$r_O = \frac{2-\sqrt3}{\sqrt3}R = 0.155\,R$ y
$r_T = \frac{\sqrt5-\sqrt3}{\sqrt3}R = 0.291\,R$. En la BCC, el \omterm{def:b1:crystals-metals:site}{hueco
tetraédrico} es el mayor.
\end{solution}

\begin{solution}{exo:b1:crystals-metals:11}
$C = Z\cdot\frac43\pi R^3/a^3$ y $Z/a^3 = \rho N_A/M$, de modo que $C =
\frac{4\pi R^3\rho N_A}{3M} = \frac{\pi\rho N_A}{6M}(2R)^3$. Cobre:
$\pi \times 8.93 \times 6.022 \times 10^{23}/(6 \times 63.55) \times (2.556
\times 10^{-8})^3 = 0.740$.
\end{solution}

\begin{solution}{exo:b1:crystals-metals:12}
$a \approx (407.8 + 408.6)/2 = \qty{408.2}{pm}$; masa molar media $(196.97 +
107.87)/2 = \qty{152.42}{g/mol}$; $\rho = 4 \times 152.42/(6.022 \times
10^{23} \times (4.082 \times 10^{-8})^3) = \qty{14.9}{g/cm^3}$.
\end{solution}

\begin{solution}{pb:b1:crystals-metals:1}
\textbf{1.} $Z = 2$; coordinación 8.
\textbf{2.} A lo largo de la diagonal principal: $a\sqrt3 = 4R$, $R = 286.65 \times
\sqrt3/4 = \qty{124.1}{pm}$.
\textbf{3.} $\rho = 2 \times 55.845/(6.022 \times 10^{23} \times (2.8665
\times 10^{-8})^3) = \qty{7.87}{g/cm^3}$.
\textbf{4.} $2/a^3 = 2/(2.8665 \times 10^{-8})^3 = \num{8.49e22}$ átomos.
\textbf{5.} $C = \pi\sqrt3/8 = 0.680$.
\textbf{6.} Las densidades calculada y medida coinciden en tres cifras.
\textbf{7.} $\alpha$: $R = 289.8\sqrt3/4 = \qty{125.5}{pm}$; $\gamma$:
$R = 363.9\sqrt2/4 = \qty{128.7}{pm}$.
\textbf{8.} $\alpha$: $a^3/2 = \qty{12.17e6}{pm^3}$ por átomo; $\gamma$:
$a^3/4 = \qty{12.05e6}{pm^3}$ por átomo.
\textbf{9.} La fase $\gamma$ es más densa: $\rho_\alpha = 55.845/(6.022 \times 10^{23}
\times 1.217 \times 10^{-23}) = \qty{7.62}{g/cm^3}$, $\rho_\gamma =
\qty{7.70}{g/cm^3}$.
\textbf{10.} $(12.05 - 12.17)/12.17 = -1.0\,\%$: la barra se contrae
alrededor de un 1\,\% al calentarla a través de la transición.
\textbf{11.} Sí: la FCC es más compacta (0.74 frente a 0.68), lo que pesa
más que el mayor radio atómico: $(0.680/0.740) \times (128.7/125.5)^3 = 0.99$.
\textbf{12.} $r_C = a\sqrt3/8 = 356.68 \times 0.2165 = \qty{77.2}{pm}$.
\textbf{13.} $r_O = 0.0670 \times 289.8 = \qty{19.4}{pm}$.
\textbf{14.} $r_T = 0.1260 \times 289.8 = \qty{36.5}{pm}$.
\textbf{15.} El carbono (\qty{77}{pm}) es el doble de grande que el mayor
hueco: hay que separar los átomos de hierro vecinos.
\textbf{16.} La deformación cuesta mucha energía, de modo que muy pocos
átomos de carbono entran en el hierro $\alpha$.
\textbf{17.} 4 \omterm{def:b1:crystals-metals:site}{huecos octaédricos} por celda de 4 átomos: uno por átomo de
hierro.
\textbf{18.} $r_T = (\sqrt{3/2} - 1) \times 128.7 = \qty{28.9}{pm}$.
\textbf{19.} Un carbono por hierro: \ce{FeC}, $w(\ce{C}) = 12.011/(55.845 +
12.011) = 17.7\,\%$.
\textbf{20.} Un carbono por cada diez hierros: $w = 12.011/(558.45 + 12.011) =
2.1\,\%$.
\textbf{21.} Los \omterm{def:b1:crystals-metals:site}{huecos octaédricos} del hierro $\gamma$ son mucho mayores
que cualquier hueco del hierro $\alpha$: el hierro caliente disuelve el
carbono, que el temple atrapa después.
\textbf{22.} $r_O = (\sqrt2 - 1) \times 128.7 = \textbf{\qty{53}{pm}}$, el
mayor hueco del hierro, y aun así más pequeño que un átomo de carbono
(\qty{77}{pm}): el carbono se disuelve en el hierro caliente, pero lo
dilata.
\end{solution}
